% !TeX root = ./main.tex

\chdsetup{
  title              = {面向交通信息认知与场景生成的多层次图增强技术研究},
  title*             = {Research on Multi-Level Graph-Enhanced Techniques for Traffic Information Cognition and Scene Generation},
  author             = {刘泽东},
  author*            = {Liu Zedong},
  speciality         = {交通信息工程及控制},
  speciality*        = {Transportation Information Engineering and Control},
  supervisor         = {崔华教授},
  supervisor*        = {Prof. Cui Hua},
  class-number       = {TP391}, % 分类号
  student-id         = {2020000000},  % 学号
  date               = {2026-12-01},  % 默认为今日
  % professional-type  = {专业学位类型},
  % professional-type* = {Professional degree type},
  % department         = {数学科学学院},  % 院系，本科生需要填写
  % secret-level       = {秘密},     % 绝密|机密|秘密|控阅，注释本行则公开
  % secret-level*      = {Secret},  % Top secret | Highly secret | Secret
  % secret-year        = {10},      % 保密/控阅期限
  % reviewer           = true,      % 声明页显示“评审专家签名”
  %
  % 数学字体
  % math-style         = GB,  % 可选：GB, TeX, ISO
  math-font          = xits,  % 可选：stix, xits, libertinus
}


% 加载宏包

% 定理类环境宏包
\usepackage{amsthm}

% 插图
\usepackage{graphicx}

% 实验参数曲线与知识图实例，图形文字沿用论文正文字体。
\usepackage{pgfplots}
\pgfplotsset{compat=1.18}
\usepgfplotslibrary{groupplots}
\usetikzlibrary{arrows.meta,positioning,calc}

% 三线表
\usepackage{booktabs}

% 表注
\usepackage{threeparttable}

% 跨页表格
\usepackage{longtable}

% 算法
\usepackage[ruled,linesnumbered]{algorithm2e}

% SI 量和单位
\usepackage{siunitx}

% 参考文献使用 BibTeX + natbib 宏包
% 顺序编码制
\usepackage[sort]{natbib}
\bibliographystyle{chdthesis-numerical}

% 著者-出版年制
% \usepackage{natbib}
% \bibliographystyle{chdthesis-authoryear}

% 本科生参考文献的著录格式
% \usepackage[sort]{natbib}
% \bibliographystyle{chdthesis-bachelor}

% 参考文献使用 BibLaTeX 宏包
% \usepackage[style=chdthesis-numeric]{biblatex}
% \usepackage[bibstyle=chdthesis-numeric,citestyle=chdthesis-inline]{biblatex}
% \usepackage[style=chdthesis-authoryear]{biblatex}
% \usepackage[style=chdthesis-bachelor]{biblatex}
% 声明 BibLaTeX 的数据库
% \addbibresource{bib/chd.bib}

% 配置图片的默认目录
\graphicspath{{figures/}}

% 数学命令
\makeatletter
\newcommand\dif{%  % 微分符号
  \mathop{}\!%
  \ifchd@math@style@TeX
    d%
  \else
    \mathrm{d}%
  \fi
}
\makeatother
\newcommand\eu{{\symup{e}}}
\newcommand\iu{{\symup{i}}}

% 用于写文档的命令
\DeclareRobustCommand\cs[1]{\texttt{\char`\\#1}}
\DeclareRobustCommand\env[1]{\texttt{#1}}
\DeclareRobustCommand\pkg[1]{\textsf{#1}}
\DeclareRobustCommand\file[1]{\nolinkurl{#1}}

% comment 宏包用于注释掉大段内容
\usepackage{comment}

\usepackage{xcolor}

\newcommand{\red}[1]{\textcolor{red}{#1}}
\newcommand{\blue}[1]{\textcolor{blue}{#1}}

% hyperref 宏包在最后调用
\usepackage{hyperref}
